\begin{statementbox}
	\textbf{\textcolor{CStatement}{条件}}
	\begin{description}[style=nextline, leftmargin=2.8em, labelsep=0.5em, itemsep=0.35em]
		\item[\textbf{\textcolor{CStatement}{条件1（二次系数为正）：}}]
			不等式必须化为标准形式 $ax^2+bx+c>0$（或 $<0$），且 $a>0$。
		\item[\textbf{\textcolor{CStatement}{条件2（判别式与根）：}}]
			记判别式 $\Delta=b^2-4ac$。当 $\Delta\ge0$ 时，方程 $ax^2+bx+c=0$ 的两根为 $x_{1,2}=\frac{-b\pm\sqrt{\Delta}}{2a}$，约定 $x_1\le x_2$。
	\end{description}
	\textbf{\textcolor{CStatement}{结论}}
	\begin{description}[style=nextline, leftmargin=2.8em, labelsep=0.5em, itemsep=0.45em]
		\item[\textbf{\textcolor{CStatement}{结论一（$\Delta>0$ 解集）：}}]
			\textbf{\textcolor{CStatement}{直观描述：}}
			大于取两边，小于取中间。
			\[
				ax^2+bx+c>0 \Longleftrightarrow x\in(-\infty,x_1)\cup(x_2,+\infty).
			\]
			\[
				ax^2+bx+c<0 \Longleftrightarrow x\in(x_1,x_2).
			\]
		\item[\textbf{\textcolor{CStatement}{结论二（$\Delta=0$ 解集）：}}]
			\textbf{\textcolor{CStatement}{直观描述：}}
			大于时去掉等根，小于时空集。
			\[
				ax^2+bx+c>0 \Longleftrightarrow x\in\mathbb{R}\setminus\lbrace -\frac{b}{2a}\rbrace.
			\]
			\[
				ax^2+bx+c<0 \Longleftrightarrow \emptyset.
			\]
		\item[\textbf{\textcolor{CStatement}{结论三（$\Delta<0$ 解集）：}}]
			\textbf{\textcolor{CStatement}{直观描述：}}
			大于恒成立，小于无解。
			\[
				ax^2+bx+c>0 \Longleftrightarrow x\in\mathbb{R}.
			\]
			\[
				ax^2+bx+c<0 \Longleftrightarrow \emptyset.
			\]
	\end{description}
	\textbf{\textcolor{CStatement}{常见变形（二次系数为负）}}
	\textbf{\textcolor{CStatement}{直观描述：}}
	负系不等式先转为正系数再套解。\par 若 $a<0$，则将原不等式两边同时乘以 $-1$（不等号方向反转），得到 $-ax^2-bx-c<0$（或 $>0$），此时二次项系数 $-a>0$，再按上述结论求解。
\end{statementbox}
